%! @@TITLE@@ — Unit @@UNITINT@@, Lesson @@LESSONID@@ — Lesson Plan
% PILOT SHAPE (2026-09-06, modeled on AP Statistics 1.4 minus AP Practice). Gradual-release plan:
% warm-up / guided notes & practice (I do two sections, we do the Guided Practice) / SPOKEN
% debrief / close and START THE HOMEWORK. There is no independent practice set in the notes —
% the homework is the individual practice. Teacher-facing: use the formal vocabulary and the
% lettered standard codes freely. Breakable boxes are `skillbox` + \boxguard; the two tabularx
% boxes are `fixedskillbox` (never splits). Mirror unit01/lesson02.
\documentclass[10pt]{article}
\usepackage{@@PREFIX@@-article}
\usepackage{@@PREFIX@@-boxes}
\usepackage{graphicx}   % -article does NOT load graphicx; needed for thumbnails/screenshots
\graphicspath{{images/}}
\newcommand{\TallMath}[1]{$\displaystyle #1\rule[-1.4em]{0pt}{3.2em}$}
@@COURSEMACROS@@\newcommand{\UnitNumberName}{Unit @@UNITINT@@: @@UNITTITLE@@ \quad}
\newcommand{\LessonNumberName}{Lesson @@LESSONID@@: @@TITLE@@}

\begin{document}
\begin{center}
    {\Large\bfseries \CourseName \\
    {\small\bfseries \UnitNumberName \LessonNumberName}}
\end{center}

\begin{tcolorbox}[colback=frost, colframe=cerulean, boxrule=0.9pt, arc=2mm,
    left=2mm, right=2mm, top=1.5mm, bottom=1.5mm]
    \textbf{Primary Objective:} TODO --- what students can do, interpret, and \emph{justify}
    by the end.
    \\[2pt]
    % CITE THE STANDARD CODES. They appear here and again on the Connections line.
    \textbf{Standards:} TODO (e.g.\ PS.DS.1a, PS.DS.1c)
    \\[2pt]
    \textbf{Lesson model:} \emph{Gradual release --- I do, we do, you do --- all of it inside
    the guided notes.} The notes deliver the instruction in two sections; the Guided Practice is
    worked together with students holding the pen; \textbf{the homework is the individual
    practice}, started in class while you circulate. The debrief is \emph{spoken}.
\end{tcolorbox}

\renewcommand{\arraystretch}{1.4}
\boxguard
\begin{skillbox}[Learning Targets \& Key Understandings]{goldbox}
% Fill in the \item text ONLY. Two tabularx cells, one itemize each.
\begin{tabularx}{\linewidth}{@{}X|X@{}}
\textbf{Learning targets (student language)} & \textbf{Key understandings (the ``why'')} \\[2pt]
\hline
\vspace{-6pt}
\begin{itemize}[leftmargin=1.1em, nosep, after=\vspace{2pt}]
  \item TODO: ``I can \dots'' target naming the formal term. \hfill \textit{CODE}
\end{itemize}
&
\vspace{-6pt}
\begin{itemize}[leftmargin=1.1em, nosep, after=\vspace{2pt}]
  \item \textbf{Target misconception:} TODO --- state it, and the rule that replaces it.
  \item TODO: the why, from the Understanding the Standards pages.
\end{itemize}
\end{tabularx}
\end{skillbox}

\boxguard[24]
\begin{skillbox}[{Vocabulary, Concepts, \& Theorems --- taught directly in the guided notes}]{greenbox}
\small
\renewcommand{\arraystretch}{1.35}
\begin{tabularx}{\linewidth}{@{} >{\bfseries}p{3.1cm} X @{}}
TODO term & TODO definition. \\
\end{tabularx}
\end{skillbox}

% fixedskillbox never splits, so the phase table stays intact on one page.
% THE PHASE TABLE IS A CONTRACT: 5 / 35 / 10 / 10 — I do ~20 of the 35, Guided Practice ~14;
% the debrief is spoken; the homework start is the second formative read.
\begin{fixedskillbox}[Lesson at a Glance --- \MeetingLength]{frost}
\renewcommand{\arraystretch}{1.25}
\begin{tabularx}{\linewidth}{@{}l c X X@{}}
\textbf{Phase} & \textbf{Min} & \textbf{Students} & \textbf{Teacher} \\
\hline
Warm-Up & 5 & TODO three spiral items, silent & Debrief item~3 aloud; leave TODO on the board all period \\
Guided Notes \& Practice & 35 & Fill the vocabulary box and notes 1--2 as each term is named;
  then \emph{TODO Guided Practice title} together, pen in their hand & \textbf{I do} notes 1--2
  (20) $\to$ \textbf{we do} the Guided Practice box (15) \\
Debrief & 10 & Pens down, papers up --- answer aloud, nothing to fill in & TODO what goes back
  on the board; the cold check \\
Close \& Assign & 10 & \textbf{Start the homework in class}, alone and silent & Announce the
  homework and its form (packet or Desmos), circulate for the second formative read, preview
  the next lesson \\
\end{tabularx}
\end{fixedskillbox}

\begin{fixedskillbox}[Warm-Up --- Activate Prior Knowledge (5 min)]{frost}
\begin{tabularx}{\linewidth}{@{}X X@{}}
\begin{minipage}[t]{\linewidth}\vspace{0pt}
    \textbf{The three items and what each seeds}
    \begin{itemize}[leftmargin=1.1em, itemsep=1pt, topsep=2pt]
      \item \textbf{1.} TODO. \textbf{Spirals} to Lesson TODO --- picked up in section~TODO.
      \item \textbf{2.} TODO. \textbf{Seeds:} TODO.
      \item \textbf{3.} TODO. \textbf{Seeds:} TODO.
    \end{itemize}
\end{minipage}
&
\begin{minipage}[t]{\linewidth}\vspace{0pt}
    \textbf{Running it}
    \begin{itemize}[leftmargin=1.1em, itemsep=1pt, topsep=2pt]
      \item \textbf{Debrief item~3 aloud} and leave TODO on the board all period.
      \item TODO: which item is the tell, and what it tells you.
    \end{itemize}
\end{minipage}
\end{tabularx}
\end{fixedskillbox}

\boxguard
\begin{skillbox}[Hook --- before anyone sits down]{frost}
On the board: TODO the claim. \textbf{Ask:} TODO. Take a hand vote and write the split on the
board. \textbf{Do not settle it.} Say only that the answer is the last thing in section~2.
\end{skillbox}

\boxguard
\begin{skillbox}[Guided Notes \& Practice --- I do, then we do (35 min)]{frost}
\textbf{This is the whole lesson.} There is no group activity, and there is no independent
practice set on this page: \textbf{the homework is the individual practice}, started in class.
\begin{multicols}{2}
\textbf{I do (about 20 min).} Two sections, each in two moves; students fill the blanks as each
idea is named, and the vocabulary box is filled \emph{as you go}, never in advance.

\emph{Section 1} TODO --- the first move (N min), then its second half \emph{TODO Title} (N min).

\textbf{Section 2 is the lesson.} TODO --- the first move (N min); its second half, \emph{TODO
Title}, is the crux (N min): TODO the computation worked live, the re-vote, the sentence to land.

\columnbreak
\textbf{We do (about 15 min).} The Guided Practice box, \emph{TODO title} --- the same moves the
homework will ask alone, in a fresh context. \textbf{Students hold the pen; you hold the
questions.} TODO which part is the crux transferred; take a hand count before anyone writes.

Circulate during Guided Practice and demand the reason, not the word:
\begin{itemize}[leftmargin=1.1em, nosep]
  \item \textbf{Part (a):} ``TODO prompt.''
  \item \textbf{Part (c), the crux transferred:} ``TODO prompt.''
\end{itemize}
While circulating, pick the papers to put up in the debrief --- TODO one right for the right
reason, one that fell into the trap. \textbf{Your formative read comes twice}: here, and again
in the last ten minutes as students start the homework.
\end{multicols}
\end{skillbox}

\boxguard[24]
\begin{skillbox}[Debrief --- whole class, spoken (10 min)]{frost}
Pens down, papers up. \textbf{This debrief is spoken} --- there is no box at the end of the
notes to fill in, so it lives entirely on the board and in the room.
\begin{multicols}{2}
\begin{enumerate}[leftmargin=1.3em, nosep, itemsep=3pt]
  \item TODO: put the crux back on the board and make a student say the sentence.
  \item TODO: the two Guided Practice papers you picked, read aloud.
  \item TODO: the second reflex to take.
\end{enumerate}
\columnbreak
\textbf{Check for understanding --- ask it cold, whole class.} \emph{``TODO question.''} Hands
down, thirty seconds of think time, then cold-call three students. Correct answer: TODO.
\emph{This replaces the exit ticket.}

\textbf{Formative read.} Sort the answers into three piles: \textbf{(a)} TODO right and
justified; \textbf{(b)} TODO right for the wrong reason; \textbf{(c)} TODO the misconception
survived. \textbf{If pile (c) runs past a third of the room, open the next lesson by re-running
section~2's second half}, and say so plainly.

\textbf{If the debrief runs short}, cut TODO --- item~1 cannot be cut. \textbf{Do not let the
debrief eat the last ten minutes.}
\end{multicols}
\end{skillbox}

\boxguard
\begin{skillbox}[Homework --- scored, started in class, due after two study halls]{goldbox}
TODO: the $\sim$6 items on \textbf{TODO context} and what each targets --- name the
\textbf{contrast pair}, the crux head-on, and the spiral item.

\textbf{Start it in the last ten minutes of the period, in class and alone.} \textbf{Scoring:}
TODO points --- 2 per item. \textbf{Item~TODO is where the misconception shows.}

\textbf{Packet or Desmos --- decided here.} The packet pages are authored and are the default;
assign a Desmos activity instead only where one genuinely carries this lesson's skills, and name
it. Say which one aloud at Close \& Assign.

\textbf{Preview of the next lesson.} TODO.

\textbf{Connections \& Big Ideas.} \textit{Standards covered:} TODO. \textit{Spirals forward
to:} TODO. \textit{Reaches back to:} TODO.
\end{skillbox}

\boxguard
\begin{skillbox}[Watch For (while circulating)]{redbox}
\begin{itemize}[leftmargin=1.2em, nosep]
  \item \textbf{TODO the target misconception} (notes~2, GP~(c), HW~TODO). Probe: \emph{``TODO.''}
  \item \textbf{TODO} (notes~1, HW~TODO). Probe: \emph{``TODO.''}
\end{itemize}
Cold-call prompts: \emph{``TODO.''}
\end{skillbox}

\boxguard
\begin{skillbox}[Close \& Assign (10 min)]{goldbox}
\textbf{Students start the homework here, in class, alone and silent --- this is the individual
practice, and the first minutes of it are your best formative read.} Hand the launch first ---
TODO items, scored, due at the start of the first class after two study halls. Circulate: \textbf{item~TODO tells you whether section~2
landed}; put students on TODO first and let the rest go home. Sort what you see into three piles
--- TODO (ready); TODO (ready, needs the language); TODO (the misconception survived) --- and open
the next lesson by re-running section~2's second half if that third pile ran past a third of the
room. Then close on what changed: \emph{``TODO.''} \textbf{Preview:} TODO.
\end{skillbox}

% ── Teacher notes — one per component, in packet order (three) ──────────────
% TEACHERNOTES: teacher prose lives HERE, never in a _key. No note for the debrief or the
% homework start — both are fully specified in their own boxes above.
\begin{teachernote}[Warm-Up]
    % TODO: pacing; which item is the tell; what to leave on the board.
\end{teachernote}

\begin{teachernote}[Guided Notes \& Practice]
    % TODO: pace it 20 / 15, then 10 / 10; which move to take a vote on before anyone writes;
    % where to spend the time (section 2); which Guided Practice part is the crux transferred
    % and what to protect when short; the reminder that the homework is the individual practice.
\end{teachernote}

\begin{teachernote}[Homework]
    % TODO: scoring; which items to start in class and which to let go home; which item to
    % read first when scoring and the three-pile sort; packet or Desmos; what to go over at
    % the start of the next lesson.
\end{teachernote}
\end{document}
